% ======================================================================
%  Section 4.5 --- Stability Considerations.  Corrected.
%
%  A. BROKEN REFERENCE. \ref{ch:chap:Fundamentals-of-Visual-Servoing}
%     carries both prefixes, ch: and chap:. Section 4.2.1 uses
%     \ref{chap:Fundamentals-of-Visual-Servoing}. This is why the PDF
%     prints "Chapter ??" here.
%
%  B. THE SECTION CONTRADICTS ITSELF. Paragraph two says the earlier
%     analysis "applies verbatim with L_R replaced by DL_R". Paragraph four
%     then says the system is time-varying and "the classical analysis
%     applies rigorously only to the frozen system at each instant". The
%     second is correct; the first must be weakened to match, or the
%     section concedes in its last paragraph what it asserted in its
%     second.
%
%  C. A MEASUREMENT THAT WAS NEVER MADE. "the rank requirement ... is
%     monitored in the experiments of Section 4.7.2 through the
%     conditioning of H_D". Nothing in the chapter reports the conditioning
%     of H_D, and the implementation does not compute it. This is the same
%     class of claim as 4.3.2's promise to characterise the admissible
%     occlusion fraction. The honest replacement is stronger than the
%     unperformed measurement: the floor of (4.5) makes the rank condition
%     non-binding by construction, so there is nothing to monitor.
%
%  D. THE HAT IN (4.15) IS ON THE WRONG OBJECT. What the control law
%     actually uses is a damped approximation of the pseudo-inverse, so it
%     is the pseudo-inverse that is estimated:
%         \widehat{(DL_R)^{+}}  (DL_R)
%     and not the pseudo-inverse of an estimated matrix. Both forms appear
%     in the literature, but only the first describes this control law.
%     This also connects to a measurement: 4.7.1 reports mu collapsing to
%     of order 1e-18, so near the desired pose the damped inverse tends to
%     the true pseudo-inverse and the condition is the right one to state.
% ======================================================================

\section{Stability Considerations}
\label{sec:robust-stability}

The introduction of the weighting modifies the closed-loop system and its
stability must be reconsidered, if only briefly.

The weighted control law~\eqref{eq:robust-lm-law} is the control law of
Chapter~\ref{ch:herpvs} applied to the weighted feature error
$\mathbf{D}\mathbf{e}$ and the weighted interaction matrix
$\mathbf{D}\mathbf{L}_{R}$. With the weights held fixed, the analysis of
Section~\ref{sec:redundant-stability} of
Chapter~\ref{chap:Fundamentals-of-Visual-Servoing} therefore carries over with
$\mathbf{L}_{R}$ replaced by $\mathbf{D}\mathbf{L}_{R}$: since the number of
features greatly exceeds six, local asymptotic stability in a neighbourhood of
the desired pose is guaranteed when the symmetric part of
%%% (A) single prefix. (B) "applies verbatim" -> "with the weights held
%%% fixed ... carries over", which is what the fourth paragraph says and
%%% what is actually true.
\begin{equation}
    \widehat{\bigl(\mathbf{D}\mathbf{L}_{R}\bigr)^{+}}\,
    \bigl(\mathbf{D}\mathbf{L}_{R}\bigr)
    \label{eq:robust-stability-matrix}
\end{equation}
is positive definite, the hat denoting the approximation of the pseudo-inverse
that the control law employs in place of the exact one. As in the unweighted
case, only local stability can be established, and the existence of local
minima of the weighted cost cannot be excluded.
%%% (D) The hat now covers the pseudo-inverse, and the sentence says what
%%% the hat means, which the original left to the reader.

It is worth noting that the approximation in question is the damped inverse
of~\eqref{eq:robust-lm-law} rather than a mis-estimated interaction matrix.
Since the damping $\mu$ decays with the error and is measured in
Section~\ref{sec:nominal} to reach of order $10^{-18}$ by convergence, the
damped inverse tends to the exact pseudo-inverse near the desired pose, which
is precisely the neighbourhood in which the condition above is asserted.
%%% Added. This ties the stability statement to a measurement and explains
%%% why the Gauss-Newton condition is the relevant one despite the chapter
%%% using a Levenberg-Marquardt law.

Two remarks specific to the weighted case are in order.

%%% (C) The unperformed measurement is replaced by the argument that makes
%%% it unnecessary.
First, the rank requirement noted in Section~\ref{sec:weighted-law} is the
binding constraint in principle: a sufficiently aggressive rejection, or an
occlusion covering the greater part of the informative image content, can
leave too few effective measurements, and because the damping is scaled by
$\mathrm{diag}(\mathbf{H}_{\mathbf{D}})$ it cannot restore a direction that
the weighting has removed. In the implementation this failure mode is
forestalled by the fixed floor $\varepsilon$ of~\eqref{eq:hessian-floor},
which bounds the damped Hessian away from singularity whatever the weights,
so the constraint does not bind in practice and the control law remains
defined by construction. What the floor cannot do is restore the information
that the rejected measurements carried; the consequence of removing too many
is a loss of accuracy rather than a divergence, and no divergence was observed
in any run reported in Section~\ref{sec:results}.
%%% If you would rather report the conditioning after all, it is one line:
%%%   std::cout << "cond(H_D) = " << Hsd.cond() << std::endl;
%%% inside the loop. That would let the claim be made as a measurement
%%% instead of an argument. As it stands the argument is sound and the
%%% measurement is not needed.

Second, the weights are recomputed at every iteration and are functions of the
current image, so the closed-loop system is time-varying and the analysis
above applies rigorously only to the frozen system at each instant. This is
the standard situation for iteratively reweighted schemes in visual
servoing~\cite{key26}. The step from the frozen analysis to the time-varying
system rests on the weights varying slowly relative to the control, which is
plausible because the image content evolves continuously along the trajectory
but is not verified here directly. What can be said is that the observable
consequence of a violation did not occur: the trajectories of
Section~\ref{sec:results} show no oscillation, no sign reversal and no
secondary minimum, and no instability attributable to the reweighting was
observed in any run.
%%% (F) "in practice the weights vary slowly" was an assertion. It is now
%%% labelled as the assumption it is, and the observable consequence is
%%% reported instead. A direct check is cheap if you want one: log
%%% || D_k - D_{k-1} ||_F against the iteration and confirm it is small
%%% compared with || D_k ||_F.

% ======================================================================
%  LABELS USED ABOVE
%    sec:redundant-stability                 Chapter 1, the redundant-case
%                                            stability analysis
%    chap:Fundamentals-of-Visual-Servoing    Chapter 1
%    sec:weighted-law                        4.2.2
%    eq:hessian-floor                        the floor added to 4.2.2
%    sec:nominal                             4.7.1
%    sec:results                             4.7
%
%  ONE THING I DID NOT CHANGE
%  "since the number of features greatly exceeds six" is correct as stated
%  -- P = 264000. Under occlusion the weighting rejects some 9.5% of them,
%  which leaves about 239000, so the inequality is not in danger and does
%  not need qualifying.
% ======================================================================
